\hypertarget{r29-aq1-aq1-forward-hnm-numerical-cap-thermodynamic-construction}{%
\section{HNM AQ1: numerical-cap thermodynamic construction\workstar}\label{r29:aq1:r29-aq1-aq1-forward-hnm-numerical-cap-thermodynamic-construction}}
% BEGIN CURRENT REGISTRY NOTE
\paragraph{Statement alias HNM-T-AQ1.}\label{stmt:AQ1} This identifier denotes the scoped mathematical claim admitted by this chapter\textquotesingle s gate. It adds no assumption and establishes no literature-priority claim.
\begin{quote}\small\textbf{Current extension: HNM-C-AQ2.} The historical limitation below remains the scope of this original result. The later, separately gated statement in Section~\ref{stmt:AQ2} records the extension: For the actual AQ1 centered full-Z3 subsequential state at both signs of |tau|<=1/100000000, the complete original-endpoint gauge-fixed space equals the invariant-local cyclic completion and reduces the physical energy generator. Centered Fourier tests including zero give H\_phys >= (alpha/16)(I-P\_Omega) and a simple vacuum in this representation. The full-GNS strengthening explicitly uses AM2's separately reviewed full-Hilbert finite gap. For the original xz Wilson in the same state, the actual 48-link/36-endpoint cover has seven incident stars, local reference energy <=98|tau| and variance >=61999/250000>1/5. The reverse proof separately sharpens that floor to 3099951/12500000. Its own model and state assumptions remain required.\end{quote}
% END CURRENT REGISTRY NOTE

\begin{quote}\small \textbf{Admission: accepted within scope.} At both signs of \(|\tau|\le 1/100000000\), actual centered full-Z3 whole-star boxes have a subsequence defining a compatible locally normal gauge-invariant stationary state on the norm closure of full bounded local algebras. The actual reference reset bound is 56|tau||R|; compact-resolvent cutoffs prove local trace-norm precompactness. The matched unbounded-onsite dynamics theorem and separate local-normality argument supply strongly continuous GNS evolution with a nonnegative self-adjoint physical energy generator. The forward physical cyclic restriction is reducing; full joint-fixed completion and its positive gap remain for AQ2.\par The typeset body below preserves the source-bound forward derivation. Its prospective language is historical; this boxed gate records the final reviewed verdict. The separately authored reverse and skeptical records remain linked. The common admission and attributed refinements are determined by the gate, not by a renamed equation.\end{quote}
Human project author: \textbf{Hruday N M (BUNZEEY)}. Independent forward execution of the frozen AQ1 contract, without current reverse or skeptic work. The compactness/dynamics/GNS/Fourier strategy has direct prior overlap with Gauvin~\citep{gauvin2026}, arXiv:2503.15539v3 Supplement A.10. The dynamics theorem is Nachtergaele--Sims arXiv:1410.8174v1 Section3/Theorem4.1 (also cross-checked by the source panel against NSY). This is a model-specific SU(2) application, not invention of that strategy; scientific priority is unverified.

\textbf{Result:} at either sign of \(|\tau|\le 10^{-8}\), centered complete-factor whole-star boxes have a subsequence defining a locally normal stationary state on the full coarse Z³ algebra. Its GNS evolution is strongly continuous, with a nonnegative self-adjoint physical energy generator and an invariant physical cyclic restriction. No \(\tau_*\) or HTW smallness is used. Whole-sequence convergence, state uniqueness, translation invariance and a physical spectral gap are not claimed in this loop.

\hypertarget{r29-aq1-actual-full-lattice-factor-model}{%
\subsection{Actual full-lattice factor model}\label{r29:aq1:r29-aq1-actual-full-lattice-factor-model}}

For any integer fine coordinates, set \(\pi(x,y,z)=(\lfloor x/4\rfloor,\lfloor y/2\rfloor,z)\). Euclidean division, including negative x,y, gives unique residues \(r=0,\ldots,3,s=0,1\). Each coarse factor owns the three positive links at its eight tail vertices, giving 24 untruncated SU(2) Haar factors. Its ten selected-strip links and fourteen free links are unchanged. The same selected coefficient triple, within the inherited ranges, repeats at every coarse site. All original endpoint gauge actions are retained, including heads outside tail blocks.

For \(\Lambda_N=[-N,N]^{3}\), the normalized finite operator is

\(\widehat H_N=\sum_{b\in\Lambda_N}h_b+\sum_{b+S\subset\Lambda_N}\phi_b\),

where \(h_b\ge I-P_b\), \(P_b\) is the actual selected-reference vacuum projection, \(S=\{0,e_x,e_y,e_z\}\), and \(\|\phi_b\|=M=7|\tau|\). The raw physical operator differs from delta Hhat\_N by the selected ground scalar, with \(\delta=\alpha/8\). Actual centered physical energy is \(K_N=\delta(\widehat H_N-\widehat E_N)\). AM2 supplies the finite full-Hilbert construction at the numerical cap, but a finite gap by itself is not the state construction below.

\hypertarget{r29-aq1-reset-energy-and-trace-norm-compactness}{%
\subsection{Reset energy and trace-norm compactness}\label{r29:aq1:r29-aq1-reset-energy-and-trace-norm-compactness}}

Let F be any fixed finite complete-factor region contained in \(\Lambda_N\). Reset the ground density on F to \(P_F=\bigotimes_{b\in F}P_b\), keeping its exterior marginal. The mixed trial state has finite reference form energy; nonnegative spectral truncations justify cancellation of unchanged unbounded exterior energy. Every incident retained star changes its expectation by at most 2M. Its anchor lies in F-S, whose size is at most 4\textbar F\textbar. The ground variational principle therefore gives

\begin{equation}
\operatorname{Tr}(\rho_{N,F}h_F)\le2M N_{F,N}\le8M|F|=56|\tau||F|=:C_F,
\quad h_F=\sum_{b\in F}h_b.\tag{HNM-AQ1.1}
\label{r29:aq1:eq1}
\end{equation}

All incoming stars are included. This is actual selected-reference energy, not an assumed Haar/pure-electric energy bound. At \(\tau=0\) it forces the product reference locally.

Each \(h_b\) is a compact-manifold elliptic Casimir sum plus bounded selected potential and its actual ground shift, hence has compact resolvent. So does the finite tensor sum \(h_F\). Its spectral projection \(Q_{F,L}=\mathbf1_{[0,L]}(h_F)\) is finite rank. For \(L>0\),

\begin{equation}
\operatorname{Tr}\rho_{N,F}(I-Q_{F,L})\le C_F/L,\qquad
\|\rho_{N,F}-Q_{F,L}\rho_{N,F}Q_{F,L}\|_1\le2\sqrt{C_F/L}.\tag{HNM-AQ1.2}
\label{r29:aq1:eq2}
\end{equation}

The second inequality follows by purifying the density, bounding the projected rank-one difference by 2sqrt(tail), and contracting under partial trace. For fixed L the compressed positive trace-at-most-one matrices form a compact finite-dimensional set. Equation(HNM-AQ1.2) makes the full family trace-norm precompact.

Diagonal extraction over the countable nested coarse cubes produces a subsequence \(N_k\) whose densities converge in trace norm on every finite F. Limits are positive and trace one. Partial-trace contractivity passes compatibility. Thus \(\omega_{\rm num}(A)=\operatorname{Tr}(\rho_F A)\) defines a normalized positive functional on the full local algebra and extends by norm continuity to its C*-closure. Every local restriction is normal. No compactness of an extensive global density, Haar replacement, or uniqueness of the selected subsequence is asserted.

\hypertarget{r29-aq1-exact-placement-in-the-unbounded-onsite-dynamics-theorem}{%
\subsection{Exact placement in the unbounded-onsite dynamics theorem}\label{r29:aq1:r29-aq1-exact-placement-in-the-unbounded-onsite-dynamics-theorem}}

Use the l1 metric on Z³ and \(F(r)=(1+r)^{-4}\). The shell at integer \(r\ge 1\) contains \(4r^{2}+2\) sites. Since this is at most \(6(r+1)^{2}\), and \(\sum_{r\ge1}(r+1)^{-2}\le1\),

\(\|F\|\le 7\).

In the convolution sum, either \(d(x,z)\ge d(x,y)/2\) or \(d(z,y)\ge d(x,y)/2\). On that half, the corresponding F factor is at most 16F(d(x,y)); summing both halves gives convolution constant \(C\le 32\|F\|\le 224\).

Whole stars have l1 diameter two. Each site belongs to four anchor stars. Therefore the actual normalized interaction satisfies

\begin{equation}
J\le4M=28|\tau|,\qquad
\|\Phi\|_F\le3^4J=81J\le2268|\tau|.\tag{HNM-AQ1.3}
\label{r29:aq1:eq3}
\end{equation}

This is a finite interaction norm, not the extensive total norm. The countable lattice, separable onsite Hilbert spaces, self-adjoint unbounded \(h_b\), and bounded self-adjoint finite-range Phi(X) now satisfy Nachtergaele--Sims Section3 and Theorem4.1. Its native restriction \(\sum_{X\subset\Lambda}\Phi(X)\) is exactly the retained whole-star prescription here. The theorem gives, for each bounded local A, the norm limit of its finite evolutions, uniform on compact \textbf{normalized} time intervals. It extends to a group of isometric star automorphisms of the quasi-local algebra.

Restore physical time by \(u=\delta t/\hbar\):

\begin{equation}
\mathcal T_t(A)=\lim_N e^{iu\widehat H_N}A e^{-iu\widehat H_N}
=\lim_N e^{itK_N/\hbar}A e^{-itK_N/\hbar}.\tag{HNM-AQ1.4}
\label{r29:aq1:eq4}
\end{equation}

All raw/reference ground scalars cancel in this conjugation. A physical compact time window corresponds to a fixed compact u window. The theorem does \textbf{not} assert norm continuity in time for every operator of the full local B(H) algebra. Its proof handles strongly continuous interaction-picture operators; no unjustified norm differentiability is added here.

\hypertarget{r29-aq1-stationarity-and-separate-gns-continuity}{%
\subsection{Stationarity and separate GNS continuity}\label{r29:aq1:r29-aq1-stationarity-and-separate-gns-continuity}}

Local density convergence extends along the same subsequence to every fixed quasi-local observable by norm approximation. Finite ground invariance and the norm volume limit imply

\(\omega_{\rm num}(\mathcal T_t(A))=\omega_{\rm num}(A)\).

Indeed replace \(T_t(A)\) by a fixed finite-region approximation, pass the local state limit, then use the uniform norm errors to compare with each finite-volume invariant expectation. The construction yields well-defined unitaries

\(U_t\pi(A)\Omega=\pi(\mathcal T_t(A))\Omega\), with \(U_t\) \(\Omega=\Omega\).

Strong time continuity requires a separate argument. For a fixed finite containing region D, its finite unitary evolution is strong-* continuous on every bounded A. The normal density \(\rho_D\) makes \(\omega_{\rm num}(A^*\mathcal T_t^D(A))\) continuous by bounded dominated convergence in its trace-class decomposition. Uniform-on-compact-time norm volume approximation then makes \(\omega_{\rm num}(A^*\mathcal T_t(A))\) continuous. Consequently

\(\|(U_t-I)\pi(A)\Omega\|^2=2\omega_{\rm num}(A^*A)-2\operatorname{Re}\omega_{\rm num}(A^*\mathcal T_t(A))\longrightarrow0\).

Local cyclic vectors are dense; unitarity extends this to the full GNS Hilbert space. Stone's theorem gives a self-adjoint physical energy \(H_{\mathrm{num}}\) such that \(U_t=\exp(itH_{\rm num}/\hbar)\) and \(H_{\mathrm{num}}\) \(\Omega=0\). There is no implicit reset of the physical clock.

\hypertarget{r29-aq1-nonnegative-spectral-support}{%
\subsection{Nonnegative spectral support}\label{r29:aq1:r29-aq1-nonnegative-spectral-support}}

For fixed bounded local A, the finite correlation

\(c_{N,A}(t)=\omega_N(A^*\mathcal T_t^N(A))=\langle A\Omega_N,e^{itK_N/\hbar}A\Omega_N\rangle\)

has a positive spectral measure supported on {[}0,infinity). Along the selected subsequence, local state and norm dynamical convergence give its limit \(c_A(t)=\langle\pi(A)\Omega,e^{itH_{\rm num}/\hbar}\pi(A)\Omega\rangle\). All correlations are bounded by \textbar\textbar A\textbar\textbar².

Take any nonnegative smooth compactly supported energy test f inside (-infinity,0). Its hbar-scaled Fourier transform is integrable. Fourier inversion and dominated convergence pass \(\int f\,d\nu_{N,A}=0\) to the actual Stone spectral measure \(\nu_A\). Such tests exhaust the negative half-line. Therefore the negative spectral projection annihilates every local cyclic vector and, by density, is zero. Thus \(H_{\mathrm{num}}\ge 0\). No finite positive gap is passed in this loop.

Finite ground densities are gauge invariant under every original finite endpoint group, so \(\omega_{\mathrm{num}}\) is gauge invariant. Each finite evolution takes a bounded local gauge-invariant operator to another such operator, and its norm limit lies in their norm-closed algebra \(A_{\mathrm{phys}}\). Both signs of time preserve \(A_{\mathrm{phys}}\); hence the physical cyclic space \(\overline{\pi(\mathcal A_{\rm phys})\Omega}\) reduces \(U_t\) and \(H_{\mathrm{num}}\). Its restriction is a nonnegative self-adjoint physical generator. Equality with the full joint fixed-vector sector and a nonzero excited physical signal remain for AQ2.

\hypertarget{r29-aq1-controls-and-limitations}{%
\subsection{Controls and limitations}\label{r29:aq1:r29-aq1-controls-and-limitations}}

The checker audits negative-coordinate ownership, actual endpoint counts, complete incoming-star incidence, exact shell/convolution estimates and physical units. The densities \(|e_n><e_n|\) show how local trace mass can escape without an energy-tightness bound. Inconsistent finite marginals are rejected.

On l2(N), \(U(t)e_j=e^{ijt}e_j\) is strongly continuous, while \(A e_j=e_{2j}\) obeys \(\|U(\pi/n)AU(\pi/n)^*-A\|=2\) for every n, witnessed by e\_n.~This preserves the distinction between norm volume convergence and full-B norm time continuity.

The prior strategy credit is explicit. No SU(3) coefficient, old \(\tau_*\) interval, HTW constant, prior orthant-state identity, whole-sequence limit, translation invariance or boundary uniqueness is imported. The state is a constructed full-lattice subsequential state at a fully numerical coupling cap. Subsequent physical-gap claims require the next reviewed loop.

Reproduce from the repository root:
\begin{verbatim}
python -B research/round29/forward/aq1/check.py \
  --output /absolute/new/output
\end{verbatim}

\subsection*{Final admitted limits and evidence}
\begin{itemize}
\item This is existence of a subsequential full-lattice state, not uniqueness of all states, whole-sequence convergence, translation invariance or all-boundary independence.
\item No old symbolic \(\tau_*\) or HTW restriction is used, but no identity with the earlier conditional orthant or bulk state is thereby established.
\item Norm convergence of each local evolution as volume grows is distinct from norm continuity in time on all B(H) local operators; only the separately proved GNS strong continuity is asserted.
\item The compactness/dynamics/GNS/Fourier strategy is already explicit in Gauvin Supplement A.10 and uses Nachtergaele--Sims/NSY. The contribution is the reviewed SU(2) selected-strip dictionary and numerical-cap application, not invention of the generic method.
\item No positive physical spectral gap, Wilson operator-domain identity, calibrated mass or continuum construction is admitted by this loop.
\end{itemize}
\repo{research/round29/advisor/aq1-gate.json}\par\repo{research/round29/reverse/aq1/report.md}\par\repo{research/round29/skeptic/aq1.md}

\emph{Sequencing note.} Prospective statements in this frozen derivation describe the moment of its admission. The final Round29 roadmap below governs subsequent work.
